\section{Results}\label{sec:results}

% Figures are produced by the notebooks and saved to report/figures/ (PDF preferred), e.g.:
%
% \begin{figure}[H]
%   \centering
%   \includegraphics[width=\textwidth]{learning_curves}
%   \caption{Learning curves (mean $\pm$ std over seeds); the dashed line is the random baseline.}
%   \label{fig:learning-curves}
% \end{figure}

\subsection{\ddqn{} vs.\ \ppo{}}\label{sec:res-ddqn-ppo}
\todo{Convergence speed, stability, final performance (frame from \cref{sec:vb-vs-pb}).}

\subsection{Learning curves}\label{sec:res-curves}
\todo{Return vs.\ environment steps, mean $\pm$ spread over seeds, baseline as a horizontal line.}

\subsection{Driving efficiency}\label{sec:res-efficiency}
\Cref{tab:main-results} is the template of the main comparison (metrics from \cref{sec:metrics}).
\begin{table}[H]
  \centering
  \caption{Evaluation results (aggregation as in \cref{sec:stats}).}
  \label{tab:main-results}
  \begin{tabular}{@{}llcccc@{}}
    \toprule
    Algorithm & Reward & Return & Mean speed & Distance & Completion \\
    \midrule
    Random  & default & -- & -- & -- & -- \\
    \ddqn{} & default & -- & -- & -- & -- \\
    \ppo{}  & default & -- & -- & -- & -- \\
    \bottomrule
  \end{tabular}
\end{table}
\todo{Fill in the table; add safety columns in \cref{sec:res-safety}.}

\subsection{Safety}\label{sec:res-safety}
\todo{Collision rate, collision-free rate, near-collision frequency.}

\subsection{Reward design}\label{sec:res-reward}
\todo{Table ``reward $\times$ behavioural metrics'' for default, A, B, C.
Key question: does a higher return correspond to better driving? Discuss the limits of reward-based safety.}

\subsection{Traffic density}\label{sec:res-density}
\paragraph{State of the art.}
Generalisation of RL agents to unseen conditions is a known weakness~\cite{cobbe2019generalization,kirk2023generalisation}.

\paragraph{Results.}
\todo{Matrix ``training density $\times$ test density'' (low/medium/high) for each algorithm; generalisation to unseen densities.}

\subsection{Hyper-parameter sensitivity}\label{sec:res-hp}
\paragraph{State of the art.}
Implementation and hyper-parameter choices strongly affect on-policy actor-critic
methods~\cite{andrychowicz2021matters}.

\paragraph{Results.}
\todo{Effect of learning rate, $\gam$, $\varepsilon$-schedule (\ddqn{}) and clipping $\epsilon$ (\ppo{}).}

\subsection{Qualitative behaviour}\label{sec:res-qualitative}
\todo{Frames from selected episodes: lane changes, overtaking, speed adaptation; analysis of failure cases.
Link to the video (see the appendix).}

\subsection{Computational efficiency}\label{sec:res-compute}
\todo{Env-steps to target return, gradient updates, parameter counts (hardware-independent);
wall-clock times only with the caveat about different machines.}
